\section{Conclusion et travaux futurs}
\label{sec:conclusion}

Nous avons porté un classifieur EWC de maintenance prédictive sur un microcontrôleur Cortex-M4 et établi,
dans un protocole apparié, une parité FP32 PC$\leftrightarrow$carte exacte en mode gelé, avec une latence
inférence + apprentissage continu $\ll$~\SI{100}{\milli\second} et une empreinte sous
\SI{256}{\kilo\byte} (Gaps~2 et~3). Sur cette même carte, nous avons \emph{mesuré} l'effondrement de la
quantification INT8 naïve de la tête binaire, \emph{isolé sa cause} par ablation émulée bit-exacte --- une
échelle par-tenseur non calibrée sur accumulateur trop étroit --- puis \emph{mesuré sa récupération} par un
kernel calibré, avec un F1 revenu à \num{0.9173} et \num{0.8995} et une parité gelée de \num{1.000} contre
l'émulateur. Le diagnostic reste émulé, la récupération ne l'est plus.

Deux axes complémentaires précisent ce résultat. Sur le \emph{moment} de la quantification, la calibration
du noyau récupère l'essentiel et un entraînement QAT n'ajoute pas de gain décisif : quand le pipeline
d'entraînement ne peut pas être modifié, une PTQ correctement calibrée suffit. Sur la \emph{profondeur}, la
frontière n'est pas fixée par les bits seuls mais par le couple profondeur--métrique : le ternaire tient en
AUROC et décroche en F1 mesuré sur carte, où la frontière redevient l'INT4. La granularité par-canal
repousse le décrochage, et le gain mémoire
annoncé sous 8 bits n'existe qu'avec un véritable bit-packing --- qui se paie environ
\SI{55}{\micro\second} de dépacking, sans menacer le budget.

\paragraph{Ce que l'INT8 apporte, et ce qu'il n'apporte pas.} Le budget système mesuré dissocie les trois
promesses habituelles. La RAM des poids est bien divisée par~4, mais la RAM totale du système est
inchangée (rapport \num{1.0002}), parce qu'à cette échelle de modèle les poids ne sont pas le poste
dominant ; et ce $\div 4$ vaut pour l'inférence, la tête flottante restant en mémoire dès que la carte
apprend. La latence \emph{augmente} d'environ la moitié en inférence et d'un facteur \num{2.4} en
apprentissage continu, la requantification --- un poste absent du chemin flottant --- consommant à elle
seule un tiers du budget d'inférence. L'énergie, enfin, ne bouge pas : l'écart mesuré
est un ordre de grandeur sous son incertitude. Sur cette cible, l'INT8 se justifie donc par la mémoire
seule, et le levier énergétique se trouve ailleurs --- dans la fréquence, où la marge de latence dont nous
disposons largement s'échange contre de l'autonomie.

\paragraph{Travaux futurs.} Le paradoxe de latence oriente vers des noyaux SIMD entiers (CMSIS-NN) pour le
poste MAC, et vers une chaîne à échelles en puissance de deux pour le poste de requantification ; c'est la
voie la plus directe pour que le gain mémoire cesse de se payer en temps. Le régime d'apprentissage appelle
en propre deux pistes : une requantification restreinte aux lignes modifiées par le pas de gradient, qui
attaque directement le facteur \num{2.4}, et une descente de gradient conduite dans le domaine entier, seule
manière de faire porter le gain mémoire sur l'apprentissage lui-même et non sur la seule inférence. Le
constat « la RAM totale ne
bouge pas » appelle par ailleurs une réplication sur un modèle où les poids \emph{dominent} l'empreinte,
régime dans lequel le $\div 4$ redeviendrait décisif. Enfin, trois cellules restent explicitement « à
mesurer » : le format \texttt{q15} sur carte, l'énergie par mise à jour continue et le profil énergétique
par phase. Nous les laissons vides plutôt que de les estimer --- la même règle qui, tout au long de cet
article, sépare ce qui est mesuré sur carte de ce qui est émulé sur PC.
